\chapter{Сборка полной главной кинетической суперметрики поля}
\label{ch:v8-full-field-kinetic-supermetric-assembly-gate}

Полный полевой носитель содержит тридцать скалярных направлений переноса и
двенадцать калибровочных однополевых направлений. На ненулевом импульсе их
главный символ имеет блочный вид
\begin{equation}
 \mathbb H_0=
 \begin{pmatrix}
  K_{\text{п}}\otimes I_4&0\\
  0&K_{\text{к}}\otimes P_{\mathrm T}
 \end{pmatrix},
 \qquad \rank\mathbb H_0=120+36=156.
 \label{eq:v8-full-field-principal-supermetric}
\end{equation}
Нулевой смешанный блок относится именно к главному символу. Скаляр и
четырёхвектор имеют разные тензорные типы. Члены с фоном, массой или одной
производной относятся к нижнему порядку и здесь не исключаются.

До фиксации всё ядро сосредоточено в продольной калибровочной части:
\begin{equation}
 \dim\ker\mathbb H_0=12.
 \label{eq:v8-full-field-principal-kernel}
\end{equation}
При ковариантной фиксации с параметром $\xi>0$ получаем
\begin{equation}
 \mathbb H_{\xi}=
 \begin{pmatrix}
  K_{\text{п}}\otimes I_4&0\\
  0&K_{\text{к}}\otimes(P_{\mathrm T}+\xi^{-1}P_{\mathrm L})
 \end{pmatrix},
 \qquad \rank\mathbb H_{\xi}=168,
 \label{eq:v8-full-field-gauge-fixed-supermetric}
\end{equation}
и точный обратный оператор
\begin{equation}
 \mathbb H_{\xi}^{-1}=
 \begin{pmatrix}
  K_{\text{п}}^{-1}\otimes I_4&0\\
  0&K_{\text{к}}^{-1}\otimes(P_{\mathrm T}+\xi P_{\mathrm L})
 \end{pmatrix}.
 \label{eq:v8-full-field-supermetric-inverse}
\end{equation}

\begin{theorem}[Полная главная суперметрика]
Следовая метрика полного $42$-мерного репера канонически поднимается до
$168$-мерного главного кинетического символа. Его нефиксированный ранг равен
$156$, ядро имеет размерность $12$, а после фиксации оператор невырожден и
обращается поблочно.
\end{theorem}

\begin{nogocriterion}[Сборка не выбирает относительный вес секторов]
Замены $K_{\text{п}}\mapsto w_{\text{п}}K_{\text{п}}$ и
$K_{\text{к}}\mapsto w_{\text{к}}K_{\text{к}}$ при положительных
независимых весах сохраняют тензорный тип, ранги и калибровочное ядро.
Следовательно, одна сборка главного символа не выводит отношение
$w_{\text{п}}/w_{\text{к}}$, нижние порядки и абсолютный масштаб времени.
\end{nogocriterion}

\begin{protocol}
\begin{enumerate}
 \item размерность скалярного блока: \textbf{$120$};
 \item ранг поперечного калибровочного блока: \textbf{$36$};
 \item ранг полной нефиксированной суперметрики: \textbf{$156$};
 \item продольное ядро: \textbf{$12$};
 \item ранг после фиксации: \textbf{$168$};
 \item смешение типов в главном символе: \textbf{нулевое};
 \item относительный вес двух секторов: \textbf{не выведен};
 \item следующий гейт: \textbf{происхождение относительного веса из родительского действия}.
\end{enumerate}
\end{protocol}

Машинный сертификат сохранён в
\path{s2t_v8_full_field_kinetic_supermetric_assembly_gate_results.json}.